% titel: Blatt 9, A8: PCl5-Dissoziation, Kc berechnen
% skript: Mortimer Kap. 14 (Gleichgewichtskonstante)

% Teil I: Werkzeuge, kein Lösungsweg.

Eine typische Gleichgewichtsaufgabe: aus einer gegebenen
Gleichgewichtskonzentration die anderen sowie $K_c$ bestimmen.

\begin{werkzeugkasten}[Gleichgewichtsrechnung]
Reaktion $PCl_5(g) \rightleftharpoons PCl_3(g) + Cl_2(g)$.
Ausgangsstoffmengen (oft in mol/L) ansetzen und die Änderung $x$ bis ins
Gleichgewicht abziehen/addieren:
\[
  \text{Start:}\ [PCl_5]_0;\quad
  \text{GG:}\ [PCl_3]=x,\ [Cl_2]=x,\ [PCl_5]=[PCl_5]_0-x .
\]
IEC-Methode (Initial, Change, Equilibrium) strukturiert genau diese
Schritte. Ist eine Gleichgewichtskonzentration bekannt, liefert sie $x$ sofort.
\textit{Skript:} Mortimer Kap. 14.
\end{werkzeugkasten}

\begin{achtung}[Darauf achten]
\begin{itemize}
  \item Im 1\,L-Gefäß ist die Stoffmenge in mol gleich der Konzentration in
        mol/L.
  \item $[PCl_5]_{GG} = [PCl_5]_0 - [PCl_3]_{GG}$, nicht $+$.
  \item $K_c = \frac{[PCl_3][Cl_2]}{[PCl_5]}$ an der Gleichgewichtslage.
\end{itemize}
\end{achtung}
